\subsection{目标项消融与配置机制}
\label{sec:ablation_results}
表\ref{tab:diagnostic_ablation}展示相同估计、准入和计费条件下的目标项消融结果。直接持有风险预算参考的净年化收益为 \diagReferenceReturn，年化波动为 \diagReferenceVol，夏普为 \diagReferenceSharpe。完整 Global RRP 对应 \globalNetReturn、\globalVolatility{} 和 \globalSharpe。相对参考路径，第二阶段优化在本样本中保持了相近的收益水平，并降低实现波动与最大回撤；月均换手由 \diagReferenceTurnover{} 变化至 \globalMonthlyTurnover，表明风险调整伴随更高交易需求。

\begin{table}[htbp]\centering
\caption{Global RRP 目标项消融与参数冻结实验}\label{tab:diagnostic_ablation}
\resizebox{\textwidth}{!}{\begin{tabular}{lrrrrr}
\toprule
实验设定 & 净年化收益 & 年化波动 & 夏普 & 最大回撤 & 月均换手\\\midrule
Global RRP & 5.98\% & 3.58\% & 1.641 & -6.80\% & 17.29\%\\
仅风险预算参考 & 5.87\% & 4.76\% & 1.222 & -7.72\% & 9.96\%\\
收益短缺惩罚置零 & 5.38\% & 3.37\% & 1.574 & -7.08\% & 11.36\%\\
方差惩罚置零 & 5.87\% & 4.76\% & 1.222 & -7.72\% & 9.95\%\\
冻结评价首年参数 & 5.96\% & 3.55\% & 1.651 & -6.58\% & 17.67\%\\
\bottomrule
\end{tabular}}
\end{table}

将收益短缺惩罚置零后，净年化收益为 \diagNoShortfallReturn，年化波动为 \diagNoShortfallVol，夏普为 \diagNoShortfallSharpe。与完整模型相比，这一实验获得了更低的波动和换手，同时收益水平下降。由此可见，在其余年度系数固定的条件下，收益短缺项维护了参考目标附近的预测收益，并在历史实现路径中体现为收益与交易需求的共同上升。

方差惩罚置零实验对应第三章给出的退化性质：参考点同时使跟踪项与收益短缺项取零，因此最终权重回到参考组合。表中的数值结果用于核对这一理论性质。该实验说明，两项惩罚围绕参考点发挥不同作用，方差项推动降低条件估计风险，收益短缺项调节这一调整的收益代价。

\begin{table}[htbp]\centering
\caption{资本跟踪与风险贡献偏离}\label{tab:diagnostic_mechanism}
\resizebox{\textwidth}{!}{\begin{tabular}{lrrr}
\toprule
实验设定 & 平均权重距离平方 & 平均最大风险预算偏离 & 平均货币权重\\\midrule
Global RRP & 0.00232 & 6.08\% & 18.02\%\\
仅风险预算参考 & 0.00000 & 0.00\% & 16.23\%\\
收益短缺惩罚置零 & 0.00244 & 6.53\% & 18.14\%\\
方差惩罚置零 & 0.00000 & 0.00\% & 16.23\%\\
冻结评价首年参数 & 0.00244 & 6.34\% & 18.10\%\\
\bottomrule
\end{tabular}}
\end{table}

表\ref{tab:diagnostic_mechanism}将资本权重距离和风险预算偏离分别列示。完整模型的平均最大风险预算偏离为 \diagPrimaryRCError，参考组合对应 \diagReferenceRCError；该列的百分数表示相对风险贡献份额的差值。指标在每次决策日使用当期估计协方差计算，再对决策日取平均。由此，本文将宽松配置带来的风险贡献变化与实际收益、波动和成本联系起来。

\figmaybe{primary_diagnostic_ablation.pdf}{目标项消融的收益与风险预算偏离}{fig:diagnostic_ablation}

冻结评价首年参数的净年化收益为 \diagFrozenReturn，夏普为 \diagFrozenSharpe，月均换手为 \diagFrozenTurnover。该实验把参数更新的作用与凸优化结构本身区分开来。年度更新与冻结规则的差异结合下文配对抽样区间评价，年度候选集合的宽度则用于说明选择规则的统计分辨能力。
\FloatBarrier

\subsection{事前风险水平可比的辅助比较}
表\ref{tab:diagnostic_matched}展示简单组合与货币 ETF 混合后的结果。风险匹配在每个决策日使用历史协方差和主模型当期目标权重完成，全部 \diagMatchDecisions{} 个配对决策中，\diagMatchFeasible{} 个满足估计方差上限。逐期预测风险、混合比例和匹配误差保存于配套CSV。

可行性检验还发现，在主模型的Ledoit--Wolf收缩协方差下，\diagLWMatchInfeasible{}个配对决策的混合集合最低方差高于主模型目标。向共同对角尺度收缩提高了货币ETF的估计风险，使集中货币持仓在该估计下也需要承担较大方差。表中辅助结果因此采用同一历史样本协方差统一衡量双方风险，同时完整保留收缩估计的可行性记录。这一现象揭示了风险模型与工具集合的相互作用：可交易的低波动资产如何进入协方差估计，会影响可实现的风险匹配。

\begin{table}[htbp]\centering
\caption{事前风险匹配后的辅助比较}\label{tab:diagnostic_matched}
\resizebox{\textwidth}{!}{\begin{tabular}{lrrrrr}
\toprule
实验设定 & 净年化收益 & 年化波动 & 夏普 & 最大回撤 & 月均换手\\\midrule
Global RRP & 6.00\% & 3.59\% & 1.644 & -6.80\% & 17.27\%\\
Equal Weight（风险匹配） & 3.24\% & 2.65\% & 1.217 & -4.53\% & 8.21\%\\
60/40 Benchmark（风险匹配） & 2.94\% & 2.76\% & 1.064 & -4.29\% & 8.94\%\\
\bottomrule
\end{tabular}}
\end{table}

风险匹配后的 Equal Weight 净年化收益为 \diagMatchedEWReturn，实现波动为 \diagMatchedEWVol，夏普为 \diagMatchedEWSharpe；风险匹配后的 60/40 Benchmark 分别为 \diagMatchedSBReturn、\diagMatchedSBVol{} 和 \diagMatchedSBSharpe。两条辅助路径按自身漂移权重和实际交易量计费。图\ref{fig:diagnostic_matched}展示同一初始净值下的扣费增长过程。

\figmaybe{primary_diagnostic_risk_matched.pdf}{主模型与事前风险匹配辅助配置的净值}{fig:diagnostic_matched}

该比较控制的是调仓时条件估计下的方差水平。实现波动的差异同时反映资产构成、协方差预测误差和持有期间的权重漂移。因此，本文将其解释为事前风险预算相近条件下的配置比较，并同时呈现实际收益与实际风险，而以配对区块抽样衡量夏普差异的不确定性。
\FloatBarrier

\subsection{成本、执行与估计敏感性}
表\ref{tab:diagnostic_robustness}将稳健性检验统一限定于当前主模型。成本变化沿用相同的目标路径，执行延迟重新核算交易前权重，均值半衰期实验则保留年度参数路径。

\begin{table}[htbp]\centering
\caption{当前 Global RRP 的实施与估计敏感性}\label{tab:diagnostic_robustness}
\resizebox{\textwidth}{!}{\begin{tabular}{lrrrr}
\toprule
检验设定 & 净年化收益 & 年化波动 & 夏普 & 最大回撤\\\midrule
单边成本0 bps & 6.01\% & 3.58\% & 1.649 & -6.79\%\\
单边成本3 bps & 5.94\% & 3.58\% & 1.631 & -6.80\%\\
单边成本10 bps & 5.78\% & 3.58\% & 1.588 & -6.82\%\\
单边成本20 bps & 5.55\% & 3.58\% & 1.527 & -6.85\%\\
均值半衰期10日 & 5.86\% & 3.57\% & 1.614 & -6.83\%\\
均值半衰期40日 & 5.96\% & 3.59\% & 1.634 & -6.86\%\\
延后一交易日执行 & 5.84\% & 3.59\% & 1.600 & -6.91\%\\
\bottomrule
\end{tabular}}
\end{table}

在单边20 bps费用下，主模型净年化收益为 \diagCostTwentyReturn。将原目标顺延一个交易日执行后，净年化收益为 \diagDelayReturn，年化波动为 \diagDelayVol，夏普为 \diagDelaySharpe，最大回撤为 \diagDelayDD。费用与执行实验把交易条件变化映射为可观察的收益风险差异。10日与40日半衰期结果进一步说明，均值估计的响应速度会改变给定年度惩罚下的配置路径。

较短半衰期提高近期收益的影响，也提高了交易需求；较长半衰期使输入调整更平缓。其经济取舍是信号响应速度与估计噪声、成本之间的平衡，与DeMiguel等关于估计误差影响配置收益的讨论相联系\cite{demiguel2009optimal}。在本次敏感性范围内，净年化收益和风险指标保持相近数量级，说明主结果对所检验的均值响应速度具有一定稳定性。

\subsection{固定子区间比较}
表\ref{tab:diagnostic_subperiods}使用预先固定的三个连续子区间，比较主模型与两条风险匹配辅助路径。区间划分覆盖不同市场阶段，2026年只计入截至样本终点的观察。每个子区间重新以首日净值为基期计算绩效，回撤使用区间内历史高点。

\begin{table}[htbp]\centering
\caption{固定子区间中的收益、夏普与回撤}\label{tab:diagnostic_subperiods}
\resizebox{\textwidth}{!}{\begin{tabular}{llrrr}
\toprule
子区间 & 配置 & 净年化收益 & 夏普 & 最大回撤\\\midrule
2018--2020 & Global RRP & 5.95\% & 1.506 & -6.80\%\\
2018--2020 & Equal Weight（风险匹配） & 3.46\% & 1.183 & -4.53\%\\
2018--2020 & 60/40 Benchmark（风险匹配） & 3.89\% & 1.265 & -4.29\%\\
2021--2023 & Global RRP & 3.82\% & 1.204 & -2.83\%\\
2021--2023 & Equal Weight（风险匹配） & 2.20\% & 0.955 & -2.19\%\\
2021--2023 & 60/40 Benchmark（风险匹配） & 1.34\% & 0.569 & -3.22\%\\
2024--2026 & Global RRP & 8.34\% & 2.199 & -2.46\%\\
2024--2026 & Equal Weight（风险匹配） & 3.98\% & 1.474 & -1.98\%\\
2024--2026 & 60/40 Benchmark（风险匹配） & 3.54\% & 1.266 & -2.32\%\\
\bottomrule
\end{tabular}}
\end{table}

子区间结果用于识别长期平均表现内部的差异，并与完整区间指标共同解释配置的时间特征。评价标准保持既定研究设定，区间表现作为稳定性证据纳入讨论。
\FloatBarrier

\subsection{配对绩效差异与统计不确定性}
表\ref{tab:diagnostic_bootstrap}报告主模型减比较对象的夏普差及21日区块抽样95\%区间；正值表示主模型样本夏普较高。图\ref{fig:diagnostic_bootstrap}同时展示5日与63日区块长度的结果，用于评价对局部依赖长度的敏感性。

\begin{table}[htbp]\centering
\caption{配对夏普差异的区块抽样区间}\label{tab:diagnostic_bootstrap}
\resizebox{\textwidth}{!}{\begin{tabular}{lrr}
\toprule
比较对象 & 主模型减比较对象的夏普差 & 95\%区间\\\midrule
仅风险预算参考 & 0.412 & [0.184, 0.608]\\
收益短缺惩罚置零 & 0.067 & [-0.043, 0.171]\\
冻结评价首年参数 & -0.010 & [-0.062, 0.042]\\
Equal Weight（风险匹配） & 0.427 & [0.034, 0.805]\\
60/40 Benchmark（风险匹配） & 0.580 & [0.100, 1.034]\\
\bottomrule
\end{tabular}}
\end{table}

\figmaybe{primary_diagnostic_bootstrap.pdf}{不同区块长度下的配对夏普差及95\%区间}{fig:diagnostic_bootstrap}

区间整体位于零上方时，对应比较在该抽样设定下支持正向差异；区间覆盖零时，样本对差异方向的分辨程度较低。本文同时报告点估计与区间，使经济幅度、统计不确定性和实施成本共同参与结论判断。所有区间均条件于当前历史规格，研究证据定位为固定资产池与既定实施规则下的探索性经验结果。
\FloatBarrier
